% Chapter 11: Anti-Patterns and Common Mistakes
\chapter{Anti-Patterns and Common Mistakes}
\label{ch:antipatterns}

Every tool has characteristic misuse patterns.
These seven anti-patterns are the most common ways developers
undermine their own productivity with Claude Code.
Each follows the same structure: symptom, root cause, and prevention.

% -------------------------------------------------------------------
\section{Context Overload}

\antipattern{Context Overload}{%
Claude ignores important rules. Instructions from CLAUDE.md are followed
inconsistently. Behavior degrades over long sessions.}{%
Hub-and-spoke architecture for shared patterns.
Conditional rules with \code{paths:} frontmatter.
Hooks for non-negotiable standards.
Keep CLAUDE.md under 500 lines.
}

\marginnote[Warning]{If your CLAUDE.md exceeds 500 lines, split it. Context is finite; instructions compete for attention.}

\textbf{Root cause:} CLAUDE.md has grown to contain every rule, convention,
and preference accumulated over months. Claude's attention dilutes
across all of it, and critical rules get the same weight as minor preferences.

% -------------------------------------------------------------------
\section{The Kitchen Sink Session}

\antipattern{Kitchen Sink Session}{%
Performance degrades midway through a session. Claude repeats itself,
forgets earlier decisions, or produces lower-quality output than at session start.}{%
Use \code{/clear} between unrelated tasks.
One session per logical task.
Use \code{/compact} if the task genuinely requires accumulated context.
}

\textbf{Root cause:} Multiple unrelated tasks in a single session.
Debugging a database query, then writing a React component, then
reviewing a PR---each task's context remains, consuming tokens
that contribute nothing to the current task.

% -------------------------------------------------------------------
\section{Over-Correcting}

\antipattern{Over-Correcting}{%
Three or more rounds of ``no, that's not what I meant'' followed by
increasingly desperate attempts at the same task.}{%
Apply the two-failure rule: after two failed corrections,
\code{/clear} and write a better initial prompt.
The cost of accumulated correction context exceeds the cost of starting fresh.
}

\textbf{Root cause:} Each correction adds noise---the original error,
your correction, Claude's acknowledgment, and its retry.
After three rounds, the context is dominated by failure patterns.
Claude's attention is drawn to the corrections rather than the actual task.

% -------------------------------------------------------------------
\section{The Permanent Prototype}

\antipattern{Permanent Prototype}{%
Code has been ``working'' for weeks but has no tests, no type hints,
no error handling. ``We'll add those later'' becomes permanent.}{%
Define explicit phase transition checklists in CLAUDE.md.
Set a date for graduation from exploration to development.
Use hooks to enforce graduated standards.
}

\textbf{Root cause:} The exploration phase (\cref{sec:phases}) has no
defined exit criteria. Without an explicit transition, the code accumulates
users and dependencies while remaining at prototype quality.

% -------------------------------------------------------------------
\section{Pipeline Direction Violations}

\antipattern{Pipeline Direction Violations}{%
Carefully edited downstream artifacts are overwritten when an upstream
script is re-run. Hours of manual refinement are lost.}{%
Document pipeline flow direction.
Mark generated files: \code{\# GENERATED --- do not edit manually}.
If you must edit generated files, also update the source.
}

\textbf{Root cause:} The developer edits a generated file (e.g., API client
code generated from an OpenAPI spec), then re-runs the code generator.
The generator regenerates from the spec, overwriting the manual edits.

% -------------------------------------------------------------------
\section{The Verification Gap}

\antipattern{Verification Gap}{%
AI-generated code is accepted without testing. Subtle bugs surface in
production weeks later. ``It looked correct'' becomes the postmortem finding.}{%
Always provide verification criteria.
Tests before or with code---never after and separately.
Hook-enforced test requirements.
}

\textbf{Root cause:} AI-generated code \emph{looks} correct.
The syntax is clean, variable names are reasonable, the logic reads well.
This appearance of correctness is precisely why verification is essential:
the bugs are subtle, not obvious.

% -------------------------------------------------------------------
\section{Big-Bang Refactoring}

\antipattern{Big-Bang Refactoring}{%
A large rewrite attempt fails partway through, leaving the codebase in
a broken state. ``I'll just rewrite the whole module'' leads to days of
wasted work and a \code{git reset --hard}.}{%
Use the incremental refactoring protocol (\cref{sec:refactoring-protocol}):
Extract $\rightarrow$ Test $\rightarrow$ Harden $\rightarrow$ Promote.
Each step independently shippable.
Commit after each step.
}

\textbf{Root cause:} The scope of the rewrite exceeds what can be held
in context. Claude loses track of the original behavior, introduces
regressions, and the developer cannot verify correctness because the
old tests were written for the old architecture.

\begin{keyinsight}
Every anti-pattern in this chapter shares a common thread:
they arise from treating Claude as infinitely capable rather than
as a powerful but bounded system. Context is finite. Attention degrades.
Verification is essential. Working within these constraints---not
ignoring them---is what separates productive Claude usage from frustrating
Claude usage.
\end{keyinsight}
